\section*{Open questions raised by this replication}

\subsection*{Q1. Warm-start QAOA at p=2 on the (4,6)-cage}
\textbf{Basis.} The paper's headline is a cold-start-from-$|+\rangle^{\otimes n}$ comparison. Warm-start QAOA (Egger et al.\ 2021; Tate et al.\ 2023) has been shown to shift the achievable expectation upward at low $p$ on other MAX-CUT families, but the paper does not test it. If warm-start QAOA$_2$ still loses to Threshold$_2$ on the same (4,6)-cage, the paper's local-view impossibility argument (Appendix~B) extends naturally. If it wins, the headline needs the qualifier ``cold-start only''.\\
\textbf{Next steps.} Reuse \texttt{code/qaoa2\_aer.py}; add an initial layer of $R_Y(2\theta_v)$ with $\theta_v$ derived from a Goemans--Williamson SDP solution (cvxpy) on PG(2,3); optimize $(\gamma_1,\beta_1,\gamma_2,\beta_2)$ with the $R_Y$ layer fixed and then jointly. Compare against the Threshold$_2$ baseline of 0.7083. All CPU.

\subsection*{Q2. ADAPT-QAOA at depth-2-equivalent}
\textbf{Basis.} Zhu et al.\ 2022 showed ADAPT-QAOA beats standard QAOA on many MAX-CUT instances at fixed circuit depth. If ADAPT-QAOA at 2 adaptively chosen mixers beats Threshold$_2$ on PG(2,3), the paper's headline is depth-2-of-the-standard-mixer-only. If not, this strengthens the paper's impossibility argument.\\
\textbf{Next steps.} Implement ADAPT-QAOA operator-pool selection (pool = single-qubit $Y_v$ mixers) on 26 qubits (Aer statevector). Cap at 2 selected mixers to match the $p{=}2$ depth. Compare best cut fraction to 0.7083 on the same PG(2,3) graph.

\subsection*{Q3. Alternative local classical algorithms (Goemans--Williamson-local extraction, Hastings chained-local)}
\textbf{Basis.} The paper cites Hastings 2019 for Modified Threshold$_2$ at $D\!\in\!\{2,3,4,5\}$ but does not compute it side-by-side with plain Threshold$_2$ at $D{=}4$. Ranking Threshold$_1$, Threshold$_2$, Modified Threshold$_2$, Hastings-linear, and GW-local-extraction on the same graphs would show whether Threshold$_2$ is the local champion or merely one of many local winners.\\
\textbf{Next steps.} Implement Hastings 2019 \S4 linear 2-step algorithm ($\sim 10$--20 lines of numpy following the paper's Appendix~C description). Run all four on Heawood, PG(2,3), and a random girth-6 4-regular graph on $\sim 200$ vertices (random lifts). Tabulate cut-fraction rank.

\subsection*{Q4. Hardware-noise interaction}
\textbf{Basis.} The paper compares noiseless QAOA$_2$ to noise-free classical. On real hardware the gap widens. Quantifying the widening on the (4,6)-cage sharpens the practical relevance; if noisy QAOA$_2$ falls below $0.5+\varepsilon$, the claim becomes ``classical wins even in the ideal quantum limit''.\\
\textbf{Next steps.} Run \texttt{code/qaoa2\_aer.py} under \texttt{AerSimulator(method='density\_matrix', noise\_model=IBM-Falcon-like)} on Heawood (14 qubits --- PG(2,3) at 26 qubits is too large for density-matrix simulation on a laptop). Sweep two-qubit-gate error $0$ to $5\!\cdot\!10^{-3}$, measure QAOA$_2$ degradation, and cross-compare with Threshold$_2 = 0.7480$ for the same graph. Free CPU.

\subsection*{Q5. Structural predictors of the classical-QAOA$_2$ gap}
\textbf{Basis.} The paper proves the gap on any girth$>5$ $D$-regular graph via a tree-limit argument but does not empirically span the feature space (spectral gap, girth-6 vs $\ge 7$, bipartite vs non-bipartite, tree-like radius-2 fraction). It is unclear whether the $\sim 0.04$ gap at D=4 is universal or peaks for the (4,6)-cage; a family where QAOA$_2$ catches up would be genuinely new.\\
\textbf{Next steps.} Sample $\sim 50$ random 4-regular graphs at $n{=}100$--$500$, filter to girth $\in\{6,7,8\}$, split by bipartiteness. Run both algorithms (QAOA$_2$ via the paper's closed-form with per-graph tree-fraction correction; Threshold$_2$ Monte Carlo). Regress gap on spectral gap, bipartiteness, girth. All CPU.
